\chapter{Типы нейронов по электрической функции}
\chaptermark{Нейроны как приборы}
\label{sec:eetypes}

В главе~\ref{sec:neurontypes} мы рассортировали нейроны по тому, какой медиатор выбрасывает их выход. Электронщик рассортировал бы их иначе: по тому, что нейрон делает со своим сигналом, --- каким прибором он работает. Главный прибор книги мы знаем с главы~\ref{sec:glutcomputer}: интегрирующий нейрон с утечкой~\eqref{eq:glutneuron}, RC-цепь с компаратором. Но в мозге есть и другие. В таблице~\ref{tab:eetypes} они собраны по четырём группам, и почти каждый уже встречался в книге.

\begin{center}
\footnotesize
\setlength{\tabcolsep}{3pt}
\renewcommand{\arraystretch}{1.15}
\begin{longtable}{>{\raggedright}p{3.1cm}>{\raggedright}p{3.3cm}>{\raggedright}p{4.7cm}>{\raggedright\arraybackslash}p{2.4cm}}
\caption{Нейроны как приборы: название, аналог в электронике, что прибор делает и где он встречается в книге.}\label{tab:eetypes}\\
\toprule
Нейрон & Прибор & Что делает & Где в книге \\
\midrule
\endfirsthead
\toprule
Нейрон & Прибор & Что делает & Где в книге \\
\midrule
\endhead
\bottomrule
\endfoot
\multicolumn{4}{l}{\emph{Фильтры и интеграторы}} \\
интегрирующий нейрон с утечкой & RC-интегратор с компаратором & складывает входы в окне $\tau$ и срабатывает у порога & гл.~\ref{sec:glutcomputer}, \eqref{eq:glutneuron} \\
детектор совпадений & И с временным окном & срабатывает, только если входы пришли почти вместе; очень малое $\tau$ & гл.~\ref{sec:glutcomputer}, \ref{sec:code}, \eqref{eq:window} \\
фазический нейрон & фильтр верхних частот, детектор фронта & отвечает на изменение входа и замолкает & гл.~\ref{sec:sensors} \\
адаптирующийся нейрон & фильтр верхних частот с автоматической регулировкой усиления & частота падает при постоянном входе & гл.~\ref{sec:sensors}, \eqref{eq:nakarushton} \\
резонатор & колебательный контур, полосовой фильтр & лучше всего отвечает на вход своей частоты & раздел~\ref{sec:resonator} \\
интегратор без утечки & интегратор на операционном усилителе & превращает скорость в положение и держит его & раздел~\ref{sec:perfectint} \\
\midrule
\multicolumn{4}{l}{\emph{Преобразователи}} \\
тонический нейрон & преобразователь напряжение---частота & частота линейно следует за входом и почти не устаёт & гл.~\ref{sec:code} \\
градуальный нейрон (без импульсов) & аналоговый усилитель, буфер & передаёт плавное напряжение на короткое расстояние & гл.~\ref{sec:sensors} \\
выпрямляющий нейрон & однополупериодный выпрямитель & частота не бывает отрицательной, $[u]_+$ & гл.~\ref{sec:glutcomputer}, \ref{sec:gaba} \\
пара «включение---выключение» & пара выпрямителей, двухпроводный код & один отвечает на «больше», другой --- на «меньше» & гл.~\ref{sec:glutcomputer}, \eqref{eq:dualrail} \\
нейрон с шунтирующим торможением & аналоговый делитель, регулировка усиления & делит вход, а не вычитает из него & гл.~\ref{sec:gaba}, \eqref{eq:shunt} \\
\midrule
\multicolumn{4}{l}{\emph{Генераторы и время}} \\
ритмоводитель & релаксационный генератор, мультивибратор & сам выдаёт импульсы с постоянной частотой & гл.~\ref{sec:serocomputer}, \ref{sec:dofa} \\
ритмоводитель с входом & генератор, управляемый напряжением & свой ритм, который вход ускоряет или замедляет & гл.~\ref{sec:chemoutput} \\
пачечный нейрон & стробируемый генератор пачек & выдаёт пачки частых импульсов --- «тире» & гл.~\ref{sec:code}, \eqref{eq:burst} \\
нейрон, привязанный к фазе & фазовый детектор, фазовая автоподстройка & срабатывает на определённой фазе входной волны & гл.~\ref{sec:sensors}, \ref{sec:chemoutput} \\
аксон с задержкой & линия задержки & задерживает сигнал на заданное время & гл.~\ref{sec:glutcomputer}, рис.~\ref{fig:itd} \\
\midrule
\multicolumn{4}{l}{\emph{Память и переключение}} \\
бистабильный нейрон & триггер Шмитта, защёлка & короткий вход переводит его в длительное состояние & раздел~\ref{sec:bistable} \\
нейрон с дендритными ветвями & мультиплексор, маленькая многослойная сеть & ветвь пропускает вход, только если есть разрешающий вход & гл.~\ref{sec:synthesis} \\
\end{longtable}
\end{center}

Одна оговорка важнее всей таблицы: это не разные клетки, а разные \emph{режимы}. Один и тот же нейрон может быть интегратором при медленном входе и детектором совпадений, когда торможение укоротило его $\tau$ (глава~\ref{sec:code}); ритмоводителем в бодрствовании и молчаливым приёмником во сне. Какой прибор получится, решают ионные каналы мембраны и широковещательные каналы, которые их подстраивают.

\section{Четыре прибора, которых в книге ещё не было}

\subsection{Резонатор}
\label{sec:resonator}

У некоторых нейронов есть медленный ток, который противодействует любому изменению напряжения: напряжение растёт --- ток его тянет вниз, падает --- тянет вверх, но с запаздыванием. Для электронщика такой ток ведёт себя как индуктивность, а вместе с ёмкостью мембраны они образуют колебательный контур. Нейрон отвечает сильнее всего на вход определённой частоты, обычно от единиц до десятка-другого герц, и слабее --- на более медленный и более быстрый~\cite{HutcheonYarom2000}. Это полосовой фильтр в одной клетке: из шумного входа он выбирает ритм своей частоты, например местный ритм главы~\ref{sec:gaba}.

\subsection{Интегратор без утечки}
\label{sec:perfectint}

Интегратор с утечкой помнит вход лишь время $\tau$ --- десятки миллисекунд. Но глазу, чтобы стоять неподвижно, нужно помнить своё положение секундами: команда «повернуть» приходит как кратковременная скорость, а держать глаз надо в новом положении. Сеть решает это той же обратной связью, что и память главы~\ref{sec:glutcomputer}. Если группа нейронов с постоянной времени $\tau$ возбуждает сама себя с весом $w$, то $\tau\,\dd r/\dd t=-r+w\,r+I$, и постоянная времени группы равна
\begin{empheq}[box=\keybox]{equation}
\tau_{\text{эфф}} \;=\; \frac{\tau}{1-w} .
\label{eq:perfectint}
\end{empheq}
При $\tau=0.1\,\text{с}$ и $w=0.995$ получается $20\,\text{с}$: почти идеальный интегратор из нейронов, каждый из которых сам по себе забывает за десятую долю секунды~\cite{Seung1996}. Цена --- точная настройка: при $w$ чуть больше единицы интегратор уходит в насыщение, при чуть меньшей --- быстро забывает. Поэтому такому интегратору нужна постоянная подстройка обучением, вроде той, что описана в главе~\ref{sec:motorlearning}.

\subsection{Бистабильный нейрон}
\label{sec:bistable}

В главах~\ref{sec:glutcomputer} и~\ref{sec:gaba} триггер собирался из группы нейронов. Но бывает триггер и в одной клетке. У некоторых нейронов есть ток, который, однажды включившись, сам поддерживает возбуждение: короткий вход переводит нейрон в длительную работу --- \emph{плато}, --- а торможение возвращает в покой. Это триггер Шмитта: у нейрона два устойчивых состояния и гистерезис между ними. Так ведут себя, например, \nt{ACET}-нейроны, управляющие мышцами, и включается это свойство широковещательным каналом: плато появляется, когда до нейрона доходит серотонин~\cite{Hounsgaard1988}. Одна и та же клетка с \nt{SERO} --- защёлка, без него --- обычный интегратор.

\subsection{Интеграторы и резонаторы}

Теория разделяет возбудимые клетки на два класса~\cite{Izhikevich2007}. \emph{Интеграторы} похожи на генератор, управляемый напряжением, который начинает с нуля: чуть выше порога нейрон работает сколь угодно медленно, и частота плавно растёт со входом. Такой нейрон складывает всё, что пришло, и хорошо передаёт частотой величину входа. \emph{Резонаторы} похожи на генератор с минимальной частотой: медленно работать они не умеют, при пороговом входе сразу выдают импульсы с конечной частотой и предпочитают вход своей частоты. Первые --- естественные приборы частотного кода главы~\ref{sec:code}, вторые --- временного.

\begin{keyblock}
\begin{proposition}[Один нейрон --- много приборов]
\label{prop:eetypes}
По электрической функции нейроны работают как интеграторы с утечкой и без неё, детекторы совпадений, фильтры верхних и полосовых частот, преобразователи напряжение---частота, выпрямители, делители, генераторы, линии задержки и триггеры. Какой прибор получится из данной клетки, задают её ионные каналы и широковещательные каналы, которые их подстраивают. Сами режимы измерены; то, насколько мозг использует такую перестройку для вычислений, --- в основном гипотеза.
\end{proposition}
\end{keyblock}

Для инженера это похоже на программируемую логическую матрицу: железо одно, а прибор задаёт конфигурация. Только конфигурацию здесь меняют не при прошивке, а на ходу, --- медленными широковещательными каналами, о которых шла речь в главах~\ref{sec:serocomputer}--\ref{sec:acet}.

\section{Задачи}

\begin{exercise}[\lvlA]
\label{ex:18:1}
\emph{Допуск интегратора без утечки.} Нейроны с $\tau=0.1$~с держат положение глаза по~\eqref{eq:perfectint} в течение $T=1$~с с ошибкой не более $5\%$ в обе стороны. (а)~Найдите допустимый диапазон $w$. (б)~$w$ --- сумма весов $K$ синапсов, каждый с независимой ошибкой $10\%$. При каком $K$ разброс суммы укладывается в допуск?
\end{exercise}

\begin{exercise}[\lvlB]
\label{ex:18:2}
\emph{Резонатор как контур.} Медленный ток раздела~\ref{sec:resonator} опишем переменной $W$: $C\,\dd V/\dd t=-g_LV-g_wW+I$, $\tau_w\,\dd W/\dd t=V-W$. Покажите, что это $RC$-цепь с параллельной ветвью $R_w=1/g_w$, $L_w=\tau_w/g_w$, и при $C/g_L=10\ms$, $\tau_w=100\ms$, $\alpha=g_w/g_L=4$ найдите частоту резонанса и выигрыш $|Z|$ в нём над $|Z(0)|$.
\end{exercise}

\begin{exercise}[\lvlB]
\label{ex:18:3}
\emph{Как интегратор начинает работать.} (а)~Нейрон $\tau\,\dd V/\dd t=-V+\mu$, $\tau=10\ms$, с порогом $\theta$ и сбросом в ноль. При каком $\mu/\theta-1$ частота равна $1$~Гц? (б)~Какую частоту даёт модель $\tau\,\dd v/\dd t=v^2+\mu$ (импульс --- уход $v$ в $+\infty$, сброс в $-\infty$) при $\mu=0.01$?
\end{exercise}

\begin{exercise}[\lvlB]
\label{ex:18:4}
\emph{Моделирование: интегратор и резонатор.} Модель задачи~\ref{ex:18:2} с $g_L=1$, $\tau_m=10\ms$, $\tau_w=100\ms$, порогом $1$ и сбросом $V\to0$ ($W$ не меняется) получает $\mu=1.5\sin2\pi ft$, $f=1$--$40$~Гц. В какой полосе частот нейрон выдаёт импульсы при $\alpha=0$ и при $\alpha=4$? Сравните с условием $1.5\,|Z(f)|>1$.
\end{exercise}

\begin{exercise}[\lvlB]
\label{ex:18:5}
\emph{Защёлка из одной клетки.} Нейрон с током плато: $\tau\,\dd r/\dd t=-r+F(wr+I)$, $F(x)=\min(1,\max(0,x))$, $\tau=10\ms$, $w=1.5$, фон $I=-0.2$. (а)~Какой наименьшей длительности импульс $I=0.3$ переводит нейрон из $r=0$ в плато? (б)~Какой наименьшей амплитуды $B$ длинный импульс $I=-0.2-B$ возвращает его в покой?
\end{exercise}

\begin{exercise}[\lvlB]
\label{ex:18:6}
\emph{Самонастройка интегратора.} При фиксации взгляда вход $I=0$, датчик скольжения даёт скорость ухода $e=\dd r/\dd t$, и $\dd w/\dd t=-\eta\,e\,r$. При $w=0.95$, $\tau=0.1$~с, $\langle r^2\rangle=1$ найдите $\eta$, при котором за $60$~с фиксаций $|1-w|$ войдёт в допуск задачи~\ref{ex:18:1}. Что сломается, если учиться и во время поворотов глаза?
\end{exercise}

\begin{exercise}[\lvlC]
\label{ex:18:7}
\emph{Зачем перестраивать прибор?} Широковещательный канал переключает $\alpha$ модели задачи~\ref{ex:18:2} между $0$ и $4$. Во сколько раз резонатор выигрывает у интегратора в отношении сигнал/шум для слабого синуса в белом шуме тока при $11$~Гц и при $1$~Гц? Придумайте задачу, где выгодно именно переключение режима.
\end{exercise}
